\section{Related Works}

Retrieval-Augmented Generation (RAG) combines document retrieval with language generation, allowing a model to answer using information retrieved from an external corpus rather than relying only on its pretrained knowledge~\cite{lewis2020rag}. This approach is suitable for product manuals because the answer to a user question often exists in the manual, but may be difficult to find through manual browsing or keyword search.

Several related works have applied RAG to technical documentation and manual-based information retrieval. Fredqvist and Sabic~\cite{fredqvist2025ragchatbot} developed a locally hosted RAG chatbot for domain-specific manuals in an ABB Robotics case study. Their system used document preprocessing, embeddings, ChromaDB, retrieval methods such as TF-IDF and BM25, and local Llama models. Their work is closely related to this project because both systems aim to reduce manual searching in technical documents and support local deployment.

Liu et al.~\cite{liu2024automotivepdfchatbots} optimized RAG techniques for automotive industry PDF chatbots using locally deployed Ollama models. Their system addressed PDF processing, retrieval optimization, context compression, reranking, BM25 retrieval, and self-RAG for technical automotive documents. This is similar to the present project because both use local models and compare retrieval configurations for document-grounded question answering. However, the present system focuses on household and office product manuals instead of automotive industry documents.

Yildirim and Samli~\cite{yildirim2025automotiveassistant} developed a RAG-based conversational assistant for the automotive sector. Their assistant supported vehicle owners and service professionals by providing real-time, context-aware information for troubleshooting and maintenance. This relates to the current project because both systems use RAG to support users in finding technical support information, although this project focuses on multiple appliance manuals rather than automotive service data.

Argnani~\cite{argnani2025firmwaremanuals} explored a domain-specific RAG chatbot for firmware manuals. The study emphasized technical correctness, traceability, and the need for users to verify answers against relevant manual sections. This is relevant to the present work because product manual chatbots must not only generate fluent answers, but also remain grounded in source documents and provide references.

Hindy and Denton~\cite{hindy2024instructionmanual} extended RAG to instruction manual understanding using a multimodal approach with contrastive learning. Their work focuses on retrieving visual information from instruction manuals, while the present project is text-based. Nevertheless, it highlights an important direction for future improvement, since many product manuals contain diagrams, tables, and images that are not fully captured by text extraction alone.

Prior studies show that RAG is a practical approach for technical documentation, product support, and manual-based question answering. The present project differs by focusing on five appliance and office-equipment manuals, comparing three local RAG pipeline variants, and adding a trained manual classifier to automatically route user queries to the correct source manual.